\chapter{The Presentation-Validation Loop}

\begin{figure}[h]
  \centering
  \begin{tikzpicture}[
    node/.style={
      draw, rounded corners=6pt,
      minimum width=2.8cm, minimum height=1cm,
      align=center, font=\bfseries
    },
    >=Stealth, thick
  ]
    \node[node] (display)  at (90:3.2cm)  {Display};
    \node[node] (desire)   at (210:3.2cm) {Desire};
    \node[node] (dopamine) at (330:3.2cm) {Dopamine};

    \draw[->] (display)  to[bend right=20] node[midway, left,  font=\small] {elicits} (desire);
    \draw[->] (desire)   to[bend right=20] node[midway, below, font=\small] {reinforces} (dopamine);
    \draw[->] (dopamine) to[bend right=20] node[midway, right, font=\small] {drives} (display);
  \end{tikzpicture}
  \caption{The Presentation-Validation (PV) loop.}
  \label{fig:pv-loop}
\end{figure}

A recurring behavioral pattern in the context of attractiveness is the \emph{Presentation-Validation loop}, which can be summarized by three stages: \textbf{Display}, \textbf{Desire}, and \textbf{Dopamine}.

\begin{definition}
The \emph{Presentation-Validation (PV) loop} is a reinforcement cycle consisting of the following stages:
\begin{enumerate}
  \item \textbf{Display.} An individual modifies or presents a stimulus $x \in S$ --- through styling, behavior, or other means --- in a way intended or likely to increase $a(o, x, t, c)$ for observers in a given context.
  \item \textbf{Desire.} The display elicits a favorable response from one or more observers: attention, approach, or other signals that $a(o, x, t, c) > 0$. This stage requires an audience.
  \item \textbf{Dopamine.} The favorable response produces positive emotional reinforcement for the displayer, increasing the likelihood of similar presentations in the future.
\end{enumerate}
\end{definition}

\section{Relation to the Attractiveness Function}

Within the framework of Chapter~2, each stage of the PV loop has a precise interpretation.
The Display stage corresponds to a modification of the stimulus $x$, thereby shifting the inputs to $a$.
The Desire stage is the observation that $a(o, x, t, c)$ is high for observers in the relevant audience.
The Dopamine stage is the internal reinforcement signal received by the displayer in response to that observation.

Over time, repeated cycles reinforce the association between specific stimulus modifications and favorable responses, driving the displayer toward presentations that maximize $a$ within a given context.

\begin{remark}
The PV loop describes a pattern, not a judgment.
The motivations behind display vary --- self-expression, confidence, attraction --- and the model is agnostic to them.
At the extreme end of the Desire stage, observer responses can include unwanted attention or harassment; these are not intended outcomes of the loop, but they fall within its scope as degenerate cases.
\end{remark}
